\section{One-sided bounds on windows}\label{sec:criterion}

This section proves the analytic criterion of the paper. An upper bound $G_+(x)\le Cx^\theta$ for a positive combination of smoothed prime sums, known for $x$ in a sparse sequence of windows $Y\le x\le Y^\lambda$, excludes zeros of $\zeta$ and of the corresponding Dirichlet $L$-functions from the half-plane $\operatorname{Re}s>\theta$. The proof is a quantitative form of Landau's theorem on Dirichlet integrals of functions of one sign~\cite{landau}.

\subsection{The smoothed prime sums}\label{sec:functions}

Let $\Lambda$ denote the von Mangoldt function and $\gamma$ Euler's constant. For $x\ge1$ put
\begin{equation}\label{G}
G(x)=\sum_{n\le x}\Lambda(n)\Big(\frac xn-1\Big)-x\log x+(1+\gamma)x ,
\end{equation}
and for a nontrivial Dirichlet character $\chi$ put
\begin{equation}\label{Gchi}
G_\chi(x)=\sum_{n\le x}\Lambda(n)\operatorname{Re}\chi(n)\Big(\frac xn-1\Big)+\operatorname{Re}\frac{L'(1,\chi)}{L(1,\chi)}\,x .
\end{equation}
For $\operatorname{Re}s>1$ and every integer $n\ge1$,
\begin{gather*}
\int_n^\infty\Big(\frac xn-1\Big)x^{-s-1}\,dx=\frac{n^{-s}}{s(s-1)},\\
\int_1^\infty x\log x\;x^{-s-1}\,dx=\frac1{(s-1)^2},\qquad
\int_1^\infty x\;x^{-s-1}\,dx=\frac1{s-1}.
\end{gather*}
Summation over $n$ gives, for $\operatorname{Re}s>1$,
\begin{equation}\label{F}
\int_1^\infty G(x)x^{-s-1}\,dx=F(s),\qquad \int_1^\infty G_\chi(x)x^{-s-1}\,dx=\tfrac12\big(F_\chi(s)+F_{\bar\chi}(s)\big),
\end{equation}
where
\begin{equation}\label{Fchi}
F(s)=-\frac{\zeta'(s)}{\zeta(s)}\,\frac1{s(s-1)}-\frac1{(s-1)^2}+\frac{1+\gamma}{s-1},\qquad
F_\chi(s)=-\frac{L'(s,\chi)}{L(s,\chi)}\,\frac1{s(s-1)}+\frac{L'(1,\chi)}{L(1,\chi)}\,\frac1{s-1}.
\end{equation}
The expansions $-\zeta'(s)/\zeta(s)=(s-1)^{-1}-\gamma+O(s-1)$ and $1/(s(s-1))=(s-1)^{-1}-1+O(s-1)$ show that $F$ is analytic at $s=1$. The same computation, together with $L(1,\chi)\ne0$, shows that $F_\chi$ is analytic at $s=1$. Both functions are meromorphic on $\mathbb C$. In the half-plane $\operatorname{Re}s>0$ their poles are simple and lie at the zeros $\rho$ of $\zeta$ and of $L(\cdot,\chi)$, with residues
\begin{equation}\label{res}
\Res_{s=\rho}F=-\frac{m_\zeta(\rho)}{\rho(\rho-1)},\qquad \Res_{s=\rho}F_\chi=-\frac{m_\chi(\rho)}{\rho(\rho-1)},
\end{equation}
where $m_\zeta(\rho)$ and $m_\chi(\rho)$ denote the orders of vanishing of $\zeta$ and of $L(\cdot,\chi)$ at $\rho$.

We use four classical facts~\cite{davenport}. The functions $\zeta$ and $L(\cdot,\chi)$ do not vanish on the half-plane $\operatorname{Re}s\ge1$, and $\zeta$ does not vanish on the interval $0<s<1$. Mertens' estimate $\sum_{n\le x}\Lambda(n)/n=\log x+O(1)$ and Chebyshev's estimate $\sum_{n\le x}\Lambda(n)=O(x)$ hold, and for nontrivial $\chi$ the partial sums $\sum_{n\le x}\Lambda(n)\chi(n)/n$ are bounded. These estimates give a constant $B$ with
\begin{equation}\label{triv}
\lvert G(x)\rvert\le Bx,\qquad \lvert G_\chi(x)\rvert\le Bx\qquad(x\ge1).
\end{equation}
The function $G$ is continuous, because the term of index $n$ in \eqref{G} vanishes at $x=n$, and away from the integers $G'(x)=\sum_{n\le x}\Lambda(n)/n-\log x+\gamma$ is bounded. Hence there is a constant $K$ with
\begin{equation}\label{lip}
\lvert G(x)-G(y)\rvert\le K\lvert x-y\rvert\qquad(x,y\ge1),
\end{equation}
and the same argument gives \eqref{lip} for $G_\chi$.

\subsection{Concentration of the kernel and the nearest poles}

\begin{lemma}\label{tails}
Let $I(t)=t-1-\log t$. For $n\ge1$ an integer, the density
$e^{-v}v^{n-1}/(n-1)!$ on $v>0$ has mass at most $e^{-nI(t)}$
on $v\le nt$ if $0<t<1$, and on $v\ge nt$ if $t>1$.
\end{lemma}

\begin{proof}
Write $V$ for a variable with this density. Substitution in the gamma
integral gives $\mathbb E e^{sV}=(1-s)^{-n}$ for $s<1$.
Take $s=1-1/t$. On either indicated tail $sV\ge snt$, so its
probability is at most $e^{-snt}(1-s)^{-n}=e^{-nI(t)}$.
\end{proof}

\begin{lemma}\label{rot}
Let $0<\varphi<\pi/2$ and $\alpha\in\mathbb R$. Every block of
$\lceil2\pi/\varphi\rceil$ consecutive integers contains a $j$
with $\cos(\alpha+j\varphi)\ge\tfrac12$.
\end{lemma}

\begin{proof}
Around the circle, the angles in the block leave gaps of length at most
$\varphi$, including the gap back to the first angle. The arc on which
cosine is at least $1/2$ has length $2\pi/3>\varphi$.
\end{proof}

\begin{lemma}\label{near}
Suppose a closed discrete set $P\subset\mathbb C$, stable under conjugation,
contains $\rho_0$ with $\operatorname{Re}\rho_0>\theta$.
For arbitrarily large real $\sigma$ the nearest points of $P$ to $\sigma$
form a single real point or a conjugate pair, at distance $r<\sigma-\theta$.
\end{lemma}

\begin{proof}
The difference $(\sigma-\theta)^2-|\sigma-\rho_0|^2$ grows linearly
with positive leading coefficient $2(\operatorname{Re}\rho_0-\theta)$.
For large $\sigma$ the nearest points therefore lie within $\sigma-\theta$.
Two distinct points of $P$ can be equidistant from a real centre for
at most one value of $\sigma$, unless they are conjugates. Discard those
exceptional centres for the countably many pairs.
\end{proof}

\subsection{The window theorem}\label{sec:win}

\begin{theorem}\label{engine}
Let $H\colon[1,\infty)\to\mathbb R$ be measurable with $\lvert H(x)\rvert\le Bx$ for a constant $B$, and suppose that the function $A(s)=\int_1^\infty H(x)x^{-s-1}\,dx$, defined for $\operatorname{Re}s>1$, extends to a meromorphic function on $\mathbb C$. Let $0\le\theta<1$. Suppose that every pole $p$ of $A$ with $\operatorname{Re}p>\theta$ is simple and satisfies $\operatorname{Re}p<1$, and that the residue of $A$ at every real pole $p>\theta$ is positive. Suppose that there are constants $\lambda>1$ and $C$ such that, for arbitrarily large $Y$, the inequality $H(x)\le Cx^\theta$ holds for every $x$ with $Y\le x\le Y^\lambda$. Then $A$ does not have a pole in the half-plane $\operatorname{Re}s>\theta$.
\end{theorem}

\begin{proof}
Suppose the set $P$ of poles contains a point to the right of $\theta$.
Since $H$ is real, the poles and residues occur in conjugate pairs
away from the real axis. Fix $1<\alpha<\beta<\lambda$ and put
$\delta=\min\{I(1/\alpha),I(\lambda/\beta)\}>0$, with $I$ from
Lemma~\ref{tails}. Choose a centre $\sigma>1$ as in Lemma~\ref{near}
and write $a=\sigma-1$. All poles have real part less than $1$, so
\[
  1<r/a<(\sigma-\theta)/(\sigma-1).
\]
We can thus fix $\sigma$ large enough that $\log(r/a)<\delta$.

\emph{A lower bound from the poles.}
For $n\ge1$, differentiating the Mellin integral at $\sigma$ gives
\[
 T_n=\frac{(-1)^{n-1}A^{(n-1)}(\sigma)}{(n-1)!}
     =\int_0^\infty H(e^u)e^{-\sigma u}\frac{u^{n-1}}{(n-1)!}\,du .
\]
The growth bound justifies differentiation under the integral.
Subtract the principal parts of the nearest poles. The remainder is
analytic on a closed disc of radius $R>r$, so Cauchy's estimate gives
\[
 r^n T_n=\sum_{|\sigma-p|=r}
   (\Res_{s=p} A)\left(\frac r{\sigma-p}\right)^n+O((r/R)^n).
\]
For a real nearest pole the sum is its positive residue. For a pair, write
$\sigma-p=re^{-i\varphi}$ for the pole in the upper half-plane and
$\Res_{s=p}A=|c|e^{i\omega}$. The sum is
$2|c|\cos(\omega+n\varphi)$, with $0<\varphi<\pi/2$.
By Lemma~\ref{rot}, in either case there is a fixed $c_0>0$ such that
$r^n T_n\ge c_0$ for a sequence of $n$ with bounded gaps.

\emph{An upper bound from a window.}
Put $L=\log Y$ and take $\alpha aL\le n\le\beta aL$. On
$L\le u\le\lambda L$ the window bound gives a contribution to $T_n$
of at most $\max(C,0)(\sigma-\theta)^{-n}$.
On the two tails use $|H(e^u)|\le Be^u$ and change variables to $v=au$.
The resulting integrals are bounded by $Ba^{-n}$ times the mass of the
density in Lemma~\ref{tails} on $v\le aL$ or $v\ge a\lambda L$.
Since $aL\le n/\alpha$ and $a\lambda L\ge n\lambda/\beta$, we obtain
\[
 r^n T_n\le \max(C,0)\left(\frac r{\sigma-\theta}\right)^n
             +2B\exp\bigl(-n(\delta-\log(r/a))\bigr).
\]
Both terms tend to zero. As $Y$ increases along the given windows,
the interval $[\alpha aL,\beta aL]$ grows in length and eventually
contains an $n$ from the lower-bound sequence, a contradiction.
\end{proof}

\begin{corollary}\label{corA}
Let $\chi_1,\dots,\chi_N$ be nontrivial Dirichlet characters, let $w_1,\dots,w_N$ be positive real numbers, and put $G_+=G+\sum_{i=1}^Nw_iG_{\chi_i}$ with $G$ and $G_\chi$ as in \eqref{G} and \eqref{Gchi}. Let $0\le\theta<1$, and suppose that there are constants $\lambda>1$ and $C$ such that, for arbitrarily large $Y$, the inequality $G_+(x)\le Cx^\theta$ holds for every $x$ with $Y\le x\le Y^\lambda$. Then $\zeta(s)\ne0$ and $L(s,\chi_i)\ne0$ for every $i$ and every $s$ with $\operatorname{Re}s>\theta$.
\end{corollary}

\begin{proof}
By \eqref{F}, \eqref{Fchi} and \eqref{triv}, the function $H=G_+$ satisfies $\lvert H(x)\rvert\le(1+\sum_iw_i)Bx$, and its transform $A=F+\frac12\sum_iw_i(F_{\chi_i}+F_{\bar\chi_i})$ is meromorphic on $\mathbb C$. By \eqref{res}, the poles of $A$ in $\operatorname{Re}s>\theta$ are simple, they lie at the zeros $\rho$ of $\zeta$, of the $L(\cdot,\chi_i)$ and of the $L(\cdot,\bar\chi_i)$, and
\[
\Res_{s=\rho}A=-\frac{c_\rho}{\rho(\rho-1)},\qquad c_\rho=m_\zeta(\rho)+\tfrac12\sum_{i=1}^Nw_i\big(m_{\chi_i}(\rho)+m_{\bar\chi_i}(\rho)\big)\ge0 .
\]
These zeros satisfy $\operatorname{Re}\rho<1$. At a real zero $\rho\in(\theta,1)$ the residue is positive because $\rho(\rho-1)<0$. Theorem~\ref{engine} shows that $A$ does not have a pole with $\operatorname{Re}s>\theta$, so $c_\rho=0$ for every $\rho$ in that half-plane. Each term of $c_\rho$ is nonnegative and each $w_i$ is positive, so each multiplicity vanishes.
\end{proof}

If $E$ is an abelian number field with group $X$ of primitive Dirichlet characters, then $\zeta_E(s)=\prod_{\chi\in X}L(s,\chi)$, and the weights $w_\chi=1$ for the nontrivial $\chi\in X$ make $G_+$ the smoothed sum over the prime ideals of $E$. Part~II proves a bound of this kind for a positive combination of characters attached to one CM field, and Corollary~\ref{corA} then treats all of these characters at once.
